\section{L' équation de Dirac}
\subsection{La forme originale}
Historiquement, afin d’éviter le problème de l’énergie négative, P.A.M. Dirac a
cherché une équation quantique relativiste linéaire, comme alternative à
l’équation de Klein-Gordon,
\[
\tag{1} \boxed{(\dm\partial^\mu+\frac{m^2c^2}{\hbar^2})\Phi(x^\nu)=0
} \qquad \Rightarrow\ \partial_t\rho+\mv{\nabla}.\mv{j}=0
\]
%\[(\Box + m^2)\,\phi = 0 \]
%\[\left( \frac{1}{c^2}\frac{\partial^2}{\partial t^2}
%- \nabla^2+ \frac{m^2 c^2}{\hbar^2} \right) \phi(\vec{x}, t) = 0\]
où $\Phi$ est la fonction d'onde et $\rho$ la densité de probabilité. La
dépendance linéaire recherchée par rapport à $\partial_t$, associée à la forme
de covariance (forme $\partial^\mu$ ), conduit également à une linéarité des
dérivées spatiales $(\mv{P}=-i\hbar\mv{\nabla})$.

La forme générale est la suivante :
\[
\tag{2} \boxed{\wt{H}\ \Psi(t)={i}\;{\hbar}\;\partial_t\;\Psi(\mv{x},t)
} \qquad \text{équation de Schrödinger}
\]
avec
\[
\boxed{\wt{H}=\mv{\alpha}.\wt{\mv{P}}+\beta m} \qquad%  \{} \
\]
En particulier, il n’y a pas de termes constants pour satisfaire les conditions
(équations
linéaires en $\partial_t$) :
\[
\tag{3}
\begin{array}{ r c l }
 \wt{H}^2\Psi & = & (\wt{\mv{P}}c^2+m^2c^4)\Psi \\
 \Rightarrow\ \ E^2 & = & (\mv{P}c^2+m^2c^4) \qquad \text{énergie relativiste}
\end{array}
\]
Remarque : Il en résulte une cohérence globale de la théorie, car l’équation de
Schrödinger, associée à l’expression relativiste de l’énergie, conduit à
l’équation de Klein-Gordon (pour une particule sans spin, s = 0) ou à
l’équation de Dirac (s = 1/2).

Détermination de $\alpha$ et $\beta$ : D’après l’équation (3), on a :
\[
H^2=(\alpha_1P_1\alpha_2P_2\alpha_3P_3+\beta m)
(\alpha_1P_1\alpha_2P_2\alpha_3P_3+\beta m)
\]
\[
=\sum_{i=1}^3(\alpha_iP_i)^2+\sum_{i\neq j,\ i,j\in\{1,2,3\}}\left[
\alpha_iP_i.\alpha_jP_j+\alpha_jP_j.\alpha_iP_i\right]+
\sum_{i=1}^3\left[
\alpha_iP_i\beta+\beta\alpha_iP_i\right]m+(\beta m)^2
\]
\[
=\sum_{i=1}^3(\alpha_iP_i)^2+\sum_{i\neq j,\ i,j\in\{1,2,3\}}
(\alpha_i.\alpha_j+\alpha_j.\alpha_i)P_iP_j+
\sum_{i=1}^3\left[
\alpha_i\beta+\beta\alpha_i\right]P_im+\beta^2m^2
\]
et donc, $\{\alpha_i,\alpha_j\}=\{\alpha_i,\beta\}=0$
et $\alpha_i^2=\beta^2=1$

Quatre propriétés en découlent :

\begin{itemize}[leftmargin=1cm,  itemsep=1pt]
\item $H^\dagger=\wt{P}_i\dagger\alpha_i^\dagger+m\beta^\dagger
=H=\wt{P}_i\alpha_i+\beta m$
\item $\beta\ket{\beta}=b\ket{\beta}, donc\ \beta^2\ket{\beta}=
b(\beta\ket{\beta})\Leftrightarrow\beta^2\ket{\beta}=
b^2\ket{\beta}\Leftrightarrow b=\pm 1$
\item Conditions d’anticommutativité.
\item Tr[$\alpha_i$]=Tr[$\beta$]=0
\end{itemize}
Exemple : on peut considérer $\Psi(\mv{x},t)=
N\ph{\mv{x}}{\mv{p}(t)}\otimes\ket{\eta_{spin}}$\\

Les matrices de dimension la plus faible satisfaisant ces propriétés sont des
matrices 4 × 4.

$\Rightarrow$ Structure décrivant une particule de spin 1/2 (et son
antiparticule) :
\setlength{\fboxsep}{3pt}
\[
\boxed{
\begin{array}{  c  }
 s=\frac{1}{2},\ -\frac{1}{2}\leqslant s_z\leqslant\frac{1}{2}\\
\end{array}}
\]
\setlength{\fboxsep}{7pt}
\[
\begin{array}{ r c l }
 \wt{\mv{S}}{}^2\ket{s,s_z} & = & s(s+1)\hbar^2\ket{s,s_z}=\frac{3}{2}
\hbar^2\ket{\frac{1}{2},\pm\frac{1}{2}} \\
 \wt{S}_z\ket{s,s_z} & = & s_z\hbar\;\ket{s,s_z}=\pm\frac{\hbar}{2}
\ket{\frac{1}{2},\pm\frac{1}{2}} \\
 \wt{S}_i & = & \frac{\sigma_i\hbar}{2}\ \text{(Pauli matrices)}\\
\end{array}
\]
Les matrices de Pauli
\[
\left\{
\begin{array}{ c c c }
 \sigma_x&=&\left(\begin{array}{cc}0&1\\0&0\\\end{array}\right)\\
 \sigma_y&=&\left(\begin{array}{cc}0&-i\\i&0\\\end{array}\right)\\
 \sigma_z&=&\left(\begin{array}{cc}1&0\\0&-1\\\end{array}\right)\\
\end{array}
\right .
\]
Le choix des matrices ($\mv{\alpha}$, $\beta$) n’est pas unique.

Exemples :
\begin{itemize}[leftmargin=1cm,  itemsep=1pt]
\item La représentation de Dirac-Pauli : \[\alpha_i=\left(\begin{array}{cc}
0&\sigma_i\\\sigma_i&0\\\end{array}\right),\ \beta=\left(\begin{array}{cc}
1&0\\0&-1\\\end{array}\right)\]
\item La représentation de Weyl : \[\alpha_i=\left(\begin{array}{cc}
-\sigma_i&0\\0&\sigma_i\\\end{array}\right),\ \beta=\left(\begin{array}{cc}
0&1\\1&0\\\end{array}\right)\]
\end{itemize}
Les résultats physiques sont indépendants de ce référentiel conventionnel.
\subsection{La forme originale}
La forme générale $\beta$ :
\[\tag{4}\begin{array}{ r c l }
\beta.(\alpha_i\wt{P}_ic+\beta mc^2)\psi & = & i\hbar c\partial_t\beta\psi(x^\nu) \\
\left[\beta\alpha_i\left(-i\hbar\frac{\partial}{\partial x^i}\right)c+m\frac{c^2}{\hbar}
-i\hbar c\beta\frac{\partial}{c\partial t}\right)]\psi & = & 0 \\
 \left(i\gm\frac{\partial}{\partial x^\mu}-\frac{m}{\hbar c}\right)\psi & = & 0 \\
\end{array}\]
où $\gamma^{\mu} = \{ \beta, \beta \alpha^i \}$. On montrera que les matrices
de Dirac $\gamma^{\mu}$ sont bien des vecteurs à 4 dimensions.
\[\boxed{\sum_{k=1}^4\left\{
i\gm_{jk}\partial-m\delta_{jk}
\right\}\psi_k(x^\nu)=0}\]
Propriétés des matrices de Dirac :
\[
\tag{5} \boxed{\gm\gamma^\nu+\gamma^\nu\gm = 2g^{\mu\nu}}
\]
Exercice : Démontrez cette propriété.
\begin{itemize}
 \item $\mu=\nu=0:(\gamma^0)^2+(\gamma^0)^2=2.I$ vrai si $\beta^2=I$
 \item $\mu=1,\nu=0:\gamma^1\gamma^0+\gamma^0\gamma^1=0,
  \beta(\alpha^1\beta)+\beta(\beta\alpha^1)=0$
 \item $\mu=2,\nu=3:(\beta\alpha^2)(\beta\alpha^3)+(\beta\alpha^3)(\beta\alpha^2)=
  0, -\alpha^2\alpha^3-\alpha^3\alpha^2=0$
 \item $\mu=\nu=3:\beta\alpha^3.\beta\alpha^3+\beta\alpha^3.\beta\alpha^3=2(-1)$
\end{itemize}
On a également que
\[
\tag{6} \boxed{\gamma^{\mu\dagger} = \gamma^0\gm\gamma^0}
\]
En effet, si $\mu=0$, $(\gamma^0)^\dagger=\gamma^0\gamma^0\gamma^0$.\\
Si $\mu=i\neq 0$, $(\gamma^i)^\dagger=\gamma^0\gamma^i\gamma^0=(\gamma^0)^2\gamma^i$.\\

Il existe d'autres propriétées utiles :
\begin{itemize}
 \item $(\gamma^0)^\dagger=\gamma^0,(\gamma^0)=1$ puisque $(\gamma^0)^2=\beta$
  hermitien (voir les 4 propriétés ci-dessus);
 \item $(\gamma^i)^\dagger=-\gamma^i$ comme $(\gamma^i)^\dagger=
  (\beta\alpha^i)^\dagger=\alpha^i\beta=-\beta\alpha^i$
 \item $(\gamma^i)^2=-I$ puisque $\gamma^i\gamma^i=\beta\alpha^i\beta\alpha^i=
  -\beta^2(\alpha^i)^2=-I$
\end{itemize}
\subsection{La forme hermitienne}

Comme on pourrait prendre le conjugué complexe de l’équation de Klein-Gordon,
nous prenons ici le conjugué hermitien de l’équation (4),

\[\tag{4}\begin{array}{ r c l }
i\gamma^0\partial_0\psi+i\gamma^i\partial_i\psi-m\psi & = & 0 \\
-i\partial_0\psi^\dagger(\gamma^0)^\dagger+-i\partial_i\psi^\dagger(\gamma^i)^\dagger-
  m\psi^\dagger & = & 0 \\
-i\partial_0\psi^\dagger\gamma^0\gamma^0+-i\partial_i\psi^\dagger\gamma^i\gamma^0-
  m\psi^\dagger\gamma^0 & = & 0 \\
i(\dm\overline{\psi})\gm+m\overline{\psi} & = & 0 \\
\end{array}\]
On a donc :
\[
\tag{7} \boxed{
i\dm(\overline{\psi})\gm+m\overline{\psi}=0
}
\]
Quant à l’équation de Klein-Gordon, $\psi\times$[équation (4)]
 + [équation (7)]$\times\psi$ donne :

\[\begin{array}{ r c l }
\overline{\psi}i\gm\dm\psi-m\overline{\psi}\psi+
i(\dm\overline{\psi})\gm\psi+m\overline{\psi}\psi & = & 0+0 \\
\dm(\overline{\psi}\gm\psi) & = & 0 \\
\end{array}\]
En définissant $j^\mu\psi\gm\psi$, on a $\dm j^\mu$ (équation de
continuité). j est la densité de flux tandis que $j^0\equiv\rho$ est la
densité de probabilité (de présence spatiale).
\[
\rho=\overline{\psi}\gamma^0\psi=\psi^\dagger\gamma^0\gamma^0\psi=\sum_{k=1}^4
\psi^*_k\psi_k=\ph{\eta_s}{\eta_s}|\ph{\mv{x}}{\psi(t)}|^2\equiv\frac{dP}{d^3x}(t)>0
\]
On peut définir le courant de charge électronique : $\boxed{j^\mu =
-e\overline{\psi}\gm\psi}$.\\

L'équation (7) doit être comparée à l'équation de Klein-Gordon pour le champ
conjugué complexe $\phi^*$.
\subsection{La forme hermitienne}
Pour une particule libre, de spin 1/2 et de masse m, dans un cadre quantique relativiste, le lagrangien s'écrit
comme suit :
\[\boxed{\mc{L}=i\overline{\psi}\gm\dm\psi-m\overline{\psi}\psi}\]

$\bullet$ Équation d’Euler-Lagrange pour le champ $\psi$ :
\[
\frac{\mc{L}}{\psi}=\dm\left[\frac{\mc{L}}{\partial(\dm\psi)}\right]
\]

Exercice : Écrivez cette équation.
$-m\overline{\psi}=\dm[i\overline{\psi}\gm]\Leftrightarrow Eq.(7).$

$\bullet$ Équation d’Euler-Lagrange pour le champ $\overline{\psi}$ :
\[
\frac{\mc{L}}{\overline{\psi}}=
\dm\left[\frac{\mc{L}}{\partial(\dm\overline{\psi})}\right]
\]
Exercice : Écrivez cette équation.
$i\gm\dm\psi-m\psi=0\Leftrightarrow Eq.(4).$
Exercice : Montrez que $\mc{L}$ est hermitien.

$(\overline{\psi}\psi)^\dagger=
\psi^\dagger(\psi^\dagger\gamma^0)^\dagger=
\psi^\dagger(\gamma^0)^\dagger\psi=
\overline{\psi}\psi\ni\mc{L}$

$(\overline{\psi}i\gm\dm\psi)^\dagger=
\dm\overline{\psi}\gm\psi(-i)$,
qui donne par intégration par partie
$i\overline{\psi}\gm\dm\psi\ni\mc{L}$
\subsection{Solutions}
\subsection{Partie relative à la fonction d'onde}
Exercice : Montrez que chacune des quatre composantes du spinor $\psi_i$ satisfait
l’équation de Klein-Gordon.

Appliquons $\ga\da$ à l’équation de Dirac (4) :
\[\begin{array}{ r c l }
\ga\da i\gm\dm\psi - m\ga\da\psi & = & 0 \\
\frac{i}{2}[\ga\gm\da\dm+\gm\ga\dm\da] & = & \frac{i}{2}[\ga\gm+\gm\ga]\da\dm \\
(ig^{\alpha\mu}\da\dm+im^2)\psi & = & 0 \\
(g^{\alpha\mu}\da\dm+m^2)\psi_i & = & 0 \\
\end{array}\]
Par conséquent,
\[
\psi_j=u_je^{-\frac{i}{\hbar}p_\mu x^\mu}
\]
on peut donc identifier $u_i\equiv\ket{\eta_{spin}}$ et $\ph{x}{\psi(t)}\equiv
e^{-i/h p_\mu x^\mu}=e^{-i/h Et+i/h p.x}$.

Exercice : De quel(s) opérateur(s) $\psi_{sol}$ est-il un état propre ?

Nous avone $\wt{P}_i\psi=
-i\hbar\partial_i(ue^{-\frac{i}{\hbar}Et}e^{-\frac{i}{\hbar}\mv{p}.\mv{x}})=
-i\hbar u\; e^{-\frac{i}{\hbar}Et}\partial_ie^{-\frac{i}{\hbar}\mv{p}.\mv{x}}=
p_i\psi$, donc $\wt{P}_i\wt{P}_i\psi=p_i^2\psi$ puis $\wt{\mv{P}}\psi=\mv{p}^2\psi$
et finallement, en utilisant Eq.(3),

\[
\tag{} \boxed{
} \qquad %\{} \{} \{}  \qquad %\{} 
\]


\[
\tag{} \boxed{
} \qquad %\{} \{} \{} \[\] \qquad %\{} \[\]
\]



